\newpage
\begin{center}
\fontsize{15pt}{18pt}\selectfont 
 \bf{TÓM TẮT NỘI DUNG LUẬN VĂN}
\end{center}
\par
Luận văn nghiên cứu ứng dụng học máy trong phân loại năm nhóm nợ trên hai bộ dữ liệu của một ngân hàng thương mại tại Việt Nam, đồng thời khảo sát khả năng dự báo chuyển sang nhóm nợ cao hơn trong các tháng tiếp theo trên bộ dữ liệu có ba kỳ quan sát (hai cặp kỳ độc lập). Quy trình thực nghiệm gồm tiền xử lý, xây dựng đặc trưng, hạn chế rò rỉ thông tin, xử lý mất cân bằng, đối sánh mô hình và giải thích bằng mức độ quan trọng của đặc trưng cùng SHAP. Macro-F1 là chỉ tiêu chính của bài toán năm lớp; average precision (ký hiệu PR-AUC trong các bảng) được dùng cho bài toán chuyển nhóm.

LightGBM đạt Macro-F1 cao nhất trong các bảng so sánh chính trên tập kiểm định, bằng 0,6484 ở bộ A và 0,6930 ở bộ B. Trên kỳ dữ liệu tiếp theo của bộ B, Random Forest đạt Macro-F1 cao nhất, bằng 0,5065; Macro-F1 của các mô hình giảm tương đối từ 12,7--32,1\%. Các thí nghiệm rút gọn đặc trưng cho thấy việc bổ sung biến không luôn cải thiện kết quả. Phân tích SHAP đối chiếu giữa hai bộ dữ liệu độc lập cho kết quả nhất quán về nhóm thông tin quan trọng nhất (Top-2 trùng khớp 100\%, tương quan hạng Spearman 0,80). Với bài toán chuyển nhóm, kiểm định trên hai cặp kỳ báo cáo thật cho kết quả trái ngược: trên cặp 1 (30/04/2026--31/05/2026), XGBoost đạt average precision trung bình 0,1806 qua kiểm định chéo 5 phần, lặp 10 lần; nhưng trên cặp 2 độc lập (31/05/2026--30/06/2026), cùng cấu hình chỉ đạt 0,0053 và không mô hình nào vượt ROC-AUC 0,5521 — xấp xỉ mức ngẫu nhiên. Kiểm định chuyển giao ngoài thời gian xác nhận kết luận này. Khả năng dự báo chuyển nhóm nợ vì vậy chưa được chứng minh là ổn định theo thời gian, dù đã có bằng chứng khai thác được tín hiệu ở một giai đoạn quan sát cụ thể.

\newpage
\begin{center}
\fontsize{15pt}{18pt}\selectfont
 \bf{THESIS ABSTRACT}
\end{center}
\par
This thesis studies the application of machine learning to classify five debt groups using two datasets from a commercial bank in Vietnam, and further examines the ability to forecast transitions to a higher-risk debt group in the following months on a dataset with three reporting periods (forming two independent period pairs). The experimental pipeline includes preprocessing, feature engineering, leakage control, class-imbalance handling, model benchmarking, and interpretation using feature importance and SHAP. Macro-F1 is the primary metric for the five-class problem, while average precision (denoted PR-AUC in the tables) is used for the transition-forecasting problem.

LightGBM attains the highest Macro-F1 in the main comparison tables on the validation set, at 0.6484 on dataset A and 0.6930 on dataset B. On the subsequent reporting period of dataset B, Random Forest attains the highest Macro-F1, at 0.5065, with the models' Macro-F1 decreasing by 12.7--32.1\% in relative terms. Feature-reduction experiments show that adding more variables does not always improve results. SHAP analysis across the two independent datasets shows consistent conclusions regarding the most important information groups (100\% Top-2 overlap, Spearman rank correlation of 0.80). For the transition-forecasting problem, evaluation on two real consecutive reporting-period pairs yields contrasting results: on pair 1 (30/04/2026--31/05/2026), XGBoost achieves an average precision of 0.1806 under 5-fold cross-validation repeated 10 times; but on the independent pair 2 (31/05/2026--30/06/2026), the same configuration achieves only 0.0053, with no model exceeding a ROC-AUC of 0.5521 -- close to random chance. Out-of-time transfer validation confirms this conclusion. The ability to forecast debt-group transitions has therefore not been shown to be stable over time, even though exploitable signal was found for one specific observation period.

\newpage
\begin{center}
   \hspace{9.5cm}  Học viên thực hiện\\
 \vspace{2 cm}
   \hspace{9.5cm} \bf Đàm Công Danh
\end{center}
